\section{A current expansion with controlled contours}\label{c6:gas}

We estimate the homogeneous Laurent polynomial \(G\) in Proposition~\ref{prop:physical-pairing}, with \(m\) sites of each color. Write \(L=\pi/8,\ b=3L,\ C=2-\sqrt2\), and divide \(G\) by its coefficient at \(y^m\), namely \(C^mD(0)^{m(m-1)}\). Denote the normalized polynomial by \(P_m\), and put \(y=e^r\).

We normalize the real-line Fourier transform and its inversion by
\[
 \widehat f(p)=\int_{\mathbb R}f(s)e^{ips}\,ds,\qquad
 f(s)=\int_{\mathbb R}\widehat f(p)e^{-ips}\,\frac{dp}{2\pi}.
\]
For a function of real period \(T>0\), with \(p_n=2\pi n/T\), the same hat denotes the period integral:
\[
 \widehat f(p_n)=\int_{-T/2}^{T/2}f(s)e^{ip_ns}\,ds,\qquad
 f(s)=\frac1T\sum_{n\in\mathbb Z}\widehat f(p_n)e^{-ip_ns}.
\]
The regularized and finite-part transforms below retain this normalization. Define
\[
 \begin{gathered}
 R_s(z)=-1/(\sin(z-L)\sin(z+L)),\quad R_c(z)=\cos^2 z,\\
 e^{V(w)}=R_s(iw)R_c(iw),\quad
 \widehat\rho(p)=1/(2\cosh Lp-1)
 \end{gathered}
\]
On the real axis, take \(V(s)\) to be the real logarithm of the positive
quantity \(R_s(is)R_c(is)\); analytic branches below are continued from
this choice. Let \(\rho\) be the inverse transform of
\(\widehat\rho\), and put \(I=\int_{\mathbb R}V(s)\rho(s)\,ds\).

\begin{proposition}[Homogeneous gas upper bound]\label{prop:gas-upper}
With this normalization, as \(m\to\infty\),
\[
 \log|P_m(e^r)|\le m^2I+
 \left(\frac38+\frac{2}{3\pi^2}\Re(r^2)\right)\log m+o(\log m),
 \qquad |\Im r|<.54\pi,
\]
locally uniformly from above in \(r\).
\end{proposition}

Lowering a site label creates one or two residue particles; summing their
residues recovers the finite polynomial. The current reflection replaces
one group of these particles by inverse currents. We will first identify
the resulting sum, including its integer constraint and constant phases,
then prove a spatial energy bound for arbitrary particle configurations.
That bound is independent of \(m\); adding the external potential will
then localize the final sum.
\begin{proof}
We keep the exact phases while deriving the representation and discard
them only when estimating its absolute value.
\paragraph{Residue particles.}
Color signs \(q_g\) are \(+1,-1\); the subscript \(s\) (same), \(c\) (cross), or \(g h\) refers to these two colors.

Use spectral coordinates \(z=iw\). An elementary integration has potential \(e^{m V(w)}\), and pair weights
\[
 t_{gh}(w)=R_{gh}(iw)/[R_{gh}(iw+L)R_{gh}(iw-L)].
\]
Its measure is \(\tau dz/(2\pi i)\), \(\tau=2\sin(2L)/\sqrt C\). There are singles and doubles; a double is the residue binding two ordered particles of one color, the second at \(z-L\) relative to first coordinate \(z\) (so internally carries \(\tau\operatorname{res}_{u=-L}t_s(u/i)\) times the additional potential etc.). Other interactions multiply. Integrate centers around \(-L\) positively, divide by factorials of multiplicities of each type and color, and multiply by \(e^{(m-K)r}\) where both colors must have equal \emph{elementary} counts \(K\); sum over all nonnegative multiplicities subject to this constraint. This is the residue representation of \(P_m\). Indeed split the \(m\) poles by tiny common displacements of the \(m\) rapidities in both colors (site \(x\) contributing \(R_{gh}(z-x)\) to the potential). Between singles/doubles centered there the same-color interactions have zeros at coincidence (orders \(2,\ge1,2\) respectively). Thus each pole in each color receives at most one center. If \(A=\operatorname{res}_{z=-L}R_s(z)\), the on-site factors are \(\tau A=2/\sqrt C,\ \tau^2 A R_s(-2L)\operatorname{res}_{u=-L}t_s(u/i)=1\), corresponding to labels \(0,-1\) starting with all \(+1\). Products belonging to different sites telescope to the indicated label weights by the quotient defining \(t\) (the elementary positions bound to \(x\) are \(x-jL,\ 1\le j\le1-\alpha_x\)). Splitting can be removed by contour analyticity.

\paragraph{Periodization and mode rules.}
We periodize to put the current calculation on a finite circle, retaining an analytic nome parameter. At the end, maximum modulus in that parameter will return the bound to the original polynomial.

Use a variable \(T\) with large positive real part. Modify every sinh factor (with argument \(w+ijL\)) in \(t\) and \(e^V\) by using \(\frac12 S_T\) instead, where
\[
 S_T(w)=2\sinh w\prod_{n\ge1}(1-e^{-2nT+2w})(1-e^{-2nT-2w}).
\]
Use shifts symmetric about zero (e.g. \(\cosh^2 w\) uses \(\pm4L\)). Their multiplicities modulo 8 are respectively
\[
 t_s:(2,-1,1,0,0,0,1,-1),\quad t_c:(0,0,0,-2,2,-2,0,0),\quad
 e^V:(0,-1,0,0,2,0,0,-1).
\]
Also replace \(\tau\) by \(\tau_T\) with \(\tau_T^2=\lim_{w\to0}t_{s,T}(w)/w^2,\ \tau_\infty=\tau\), and multiply the integrand by
\(\sum_l e^{-2Tl^2+4l\sum_j q_j w_j}\).
These changes respect residues (apply before taking them). Resulting center differentials are periodic with period \(T\), by \(S_T(w+T)=-e^{2w+T}S_T(w)\). Denote the resulting Laurent polynomial by \(P_{m,T}(y)\). Its coefficients are analytic at \(e^{-2T}=0\), where they equal those of \(P_m(y)\); finiteness follows by splitting the poles as before.

First take \(T\) real. In each coefficient open the cycles into horizontal
period contours of elementary positions \(w=s+ia\). Each of the following
species occurs in both colors. The signs refer to the orientation of
the opened contour, independently of the color charge \(q_g\):
\[
\begin{array}{c|c|c}
 \text{species}&\text{contour orientation}&\text{elementary offsets}\\\hline
 R\ \text{(single)}&+&0\\
 S\ \text{(single)}&-&5L/2\\
 D\ \text{(double)}&-&31L/20,\ 51L/20
\end{array}
\]

Indeed by periodicity one starts each center on the line just below \(a=L\) minus the line just above, say at \(L\pm\varepsilon\). First move positive double centers down to \(\varepsilon\). Move the positive singles to \(0\); each crosses a same-color pole at \(w=w_{\rm neg}-iL\) (negative single at height \(L+\varepsilon\)). The residue has the opposite sign to a positive double centered at \(\varepsilon\) (top minus bottom orientation, negative partner, and the residue at \(+L\) in \(z\)-difference instead of \(-L\)). Interactions between bound particles and others cause no obstruction (same-color products \(t_s(w)t_s(w+iL),t_s(w)t_s(w-iL),t_s(w)^2t_s(w+iL)t_s(w-iL)\) have no poles, equally with the nome modification; cross poles are outside the shifts used here). Thus only disjoint pair matchings occur; dividing by factorials leaves the inverse factorials of the unpaired multiplicities and of the number of pairs. Binomial cancellation with the positive doubles removes all these bound contours. Remaining negative lines move as stated without pole crossings (zeros do not matter).

Write \(p=2\pi n/T,\ c_*=2\pi/T\). We specify mode rules carefully. For a pair \(i,j\) at offsets differing by \(A=a_i-a_j\), multiplicities \(h\) as in the arrays, put \(u=A+j'L,\ U=u\bmod\pi\in[0,\pi)\) for each shift \(j'\), and use
\[
 \begin{gathered}
 K_{ij}(p)=\frac{2\pi}{p}\sum h\left(\frac{e^{pU}}{e^{\pi p}-1}+\frac12{\bf1}_{U=0}\right),\\
 2\pi W_{ij}=\sum h(2u-2U+\pi{\bf1}_{U>0}),\quad
 k_{ij}=\sum h(U^2-\pi U+\pi^2/6).
 \end{gathered}
\]
For potentials use the last array with \(A=a_i\), writing \(V_i(p)=-K_{i,E}(p)\), \(W^V_i=W_{i,E}\). Thus
\[
 K_{ij}=4q_iq_j/p^2+(4q_iq_j A-2\pi W_{ij})/p+k_{ij}+O(p).
\]
Exact-zero arguments in the formula use an averaged branch for log absolute value; at an actual simple pole undergoing fusion one can instead take one-sided branches first (unit constant phases).

Indeed Poisson summation of the center factor gives
\(\sqrt{\pi/(2T)}\sum_l e^{2X^2/T-\pi^2 l^2/(2T)+i c_*l X}\), \(X=\sum q_i w_i\). Distribute \(2X^2\) into the pairs by neutrality. Then the log pair on each contour, up to constant phase branch, is
\[
 q_iq_j(T/3-2\log2)+(2q_iq_j A^2-k_{ij})/T+
 2\pi i W_{ij}(1/2-(s_i-s_j)/T)
 -T^{-1}\sum_{p\ne0}K_{ij}(p)e^{-ip(s_i-s_j)} .
\]
Constants here are as displayed for pairs involving \(R\) (thus no unspecified phase depending on their count). Same formula setting charge product to 0 applies to the potential. Indeed all factors are products of \(\sinh(w+iu)\) in symmetric shifts with no constant (by inspection, positive at generic real \(w\)), then of \(S_T(w+iu)/2\). On \(0<d<T\) the log of \(S_T(d+iu)\) uses leading log \(d+iu\) and periodic images of \(\log(1-e^{-2|d|-2iu\,{\rm sgn}(d)})\). Remove \((d^2+2iud)/T\) and the indicated linear phase. The remaining saw in the imaginary part, before image logs, is \(i(2U-\pi)(1/2-d/T)\) (0 instead at \(U=0\)). The means follow since the integrated image logs total \(-(U^2-\pi U+\pi^2/6)\) by the Fourier series of a quadratic. For \(p\ne0\) the negative transformed coefficient before \(1/T\) at \(0<U<\pi\) is
\[
 \frac{2}{p^2}+\frac{2U-\pi}{p}
 +\sum_{j>0}\left(\frac{e^{-2ijU}}{j(2j-ip)}+\frac{e^{2ijU}}{j(2j+ip)}\right)
 =\frac{2\pi}{p}\frac{e^{pU}}{e^{\pi p}-1},
\]
by Fourier series in \(U\), with mid-jump at 0. Combined shifts give the formula on both sides by periodicity (\(W\) integral for pairs between species here, summing both slots of a bound species); for the internal simple-pole residue use \(d\to0^+\) extracting the pole by its regular part. This latter rule can include a unit phase.

On the two \(R\) lines the matrix symbol \(K=K_{RR}\) and potential simplify to
\[
 \begin{gathered}
 K=(1-e(p))K_0,\quad e(p)=(\sinh(3Lp)-\sinh(2Lp))/\sinh(4Lp),\\
 K_0=\frac{2\pi}{p}\begin{pmatrix}\coth(bp)&-\operatorname{csch}(bp)\\-\operatorname{csch}(bp)&\coth(bp)\end{pmatrix},\quad
 V_R=K{\bf1}\widehat\rho.
 \end{gathered}
\]
In particular \(0< e(p)<1\). Put \(Q=K^{-1},\ h_*=3\pi b/4=K_+(0)\) where plus means the common-color eigenvalue.

On the preliminary contours the left block \(L'\) consists of \(S,D\)
in both colors. The inverse block has species \(I'=(g,t)\),
\(t=+,-\), shifted by \(a_+=5L/12\), \(a_-=-L/4\) from the
corresponding right lines. Introduce the larger mode matrix with blocks
\(R,L',I'\):
\[
 \begin{gathered}
 {\cal M}=\begin{pmatrix}K&K_{RL'}&-d(-p)^t\\
 K_{L'R}&K_{L'L'}&0\\
 -d(p)&0&B(p)\end{pmatrix},\quad
 d_{i}(p)= -\frac{2\pi}{p} t_i e^{a_i p}{\bf e}_{g_i}^t,\\
 B_{ij}=-\frac{2\pi}{p}\operatorname{sgn}(a_i-a_j)t_i t_j\,{\bf1}_{g_i=g_j}e^{(a_i-a_j)p}.
 \end{gathered}
\]
Use multiplication/summation over both elementary members for left doubles. Let \(J={\cal M}_{R,(L'I')}\), \(Z\) be the Schur complement removing \(K\), and 
\[
 \mu=(V_{L'}-K_{L'R}{\bf1}\widehat\rho,\ d{\bf1}\widehat\rho).
\]
Write \(n_\alpha(p)=\sum_{j\text{ of type }\alpha}e^{ip s_j}\),
and write \(\nu=n(0)\) for the vector of nonnegative integer
multiplicities. Species indices here include the color. The allowed
multiplicities satisfy
\begin{equation}\label{eq:gas-preliminary-winding}
 W_{RL'}n_{L'}+\left(\sum_{i:g_i=g}t_i n_i\right)_g=l q_R,\qquad l\in{\mathbb Z}.
\end{equation}
Here the unadorned \(n_i\) denote the counts \(n_i(0)=\nu_i\).
Equivalently total winding is zero (and \(l\) determined). Define
\(u_*=-\frac1{2h_*}\lim_{p\to0}{\bf1}^t J(p)n(p)\),
and \(L_g\) the elementary count in block \(L'\) of color \(g\); bars denote two-color averages.

Two regularizations enter the formula. At the zero Fourier mode we take
the Laurent finite part of the whole expression, including all positional
exponentials \(n_\alpha(p)\). At a particle's self-interaction we instead
take the constant remaining after subtracting the contact singularity
\(-2\log|d|\). The latter is a regular diagonal value in position space;
it does not replace the Laurent prescription at \(p=0\).

The scalar in the transform is
\begin{equation}\label{eq:gas-transform-scalar}
 {\cal P}= \sqrt{\pi/(2T)}\sqrt{\pi T/h_*}\,
 \prod_{p>0}(1-e(p))^{-2}\,
 \exp\left[\frac{m^2}{2T}\sum_p {\bf1}^t K(p){\bf1}\widehat\rho(p)^2\right],
\end{equation}
and, for a configuration of the remaining particles, write its exponent as
\begin{equation}
 \begin{aligned}
 \Phi_{m,T,r,j}(\mathbf s)&=
 \frac1T\sum_p\big(m n(-p)^t\mu(p)-\tfrac12 n(-p)^t Z(p)n(p)\big)
 \\
 &\quad+\frac{T r_j^2}{4h_*}-r_j(u_*+\bar L),
 \qquad r_j=r+2\pi i j,\quad j\in{\mathbb Z}.
 \end{aligned}
 \label{eq:legacy-P-star}
\end{equation}
The sum to be identified has the form
\begin{equation}\label{eq:gas-grand-sum}
 {\cal P}\sum_{\nu}\frac{\varepsilon_m(\nu)}{\prod_\alpha\nu_\alpha!}
 \int_{[-T/2,T/2]^{|\nu|}}
 \sum_{j\in\mathbb Z}e^{\Phi_{m,T,r,j}(\mathbf s)}
 \prod_{d=1}^{|\nu|}\frac{ds_d}{2\pi}.
\end{equation}
On the preliminary contours, \(\nu\) ranges over the multiplicities in
\(L'\cup I'\) satisfying \eqref{eq:gas-preliminary-winding}.
The unit factor \(\varepsilon_m(\nu)\) is the product of the
contour-orientation signs and constant phases from the logarithms,
current reflection, and residues. It is fixed by the derivation, rather
than replaced by its absolute value. We will show that it is independent
of the positions, \(T\), \(r\), and the integer summed before Poisson
summation, so it also lies outside the \(j\)-sum. Dependence on \(m\)
is retained.

At this stage the identity is read coefficientwise in \(y=e^r\), with
the original left counts and \(l\) fixed. We first prove absolute
convergence of that inverse-current expansion. Moving the remaining
\(D\) contours then changes the species and the factors
\(\varepsilon_m(\nu)\); the resulting full sum will be proved
absolutely convergent before it is identified with \(P_{m,T}(e^r)\).

\subsection{Current reflection and the zero modes}
We now derive the transformed sum. The nonzero modes determine a Schur complement, while the integer charge constraint determines the remaining Gaussian sum. Both are needed to fix the normalization.

\paragraph{Oscillators and the current normalization.}
For the nonzero modes use a two-component real periodic centered Gaussian field with covariance \(T^{-1}\sum_{p\ne0}(K-K_0)(p)(-1)e^{-ips}\). This implements the change from right kernel \(K\) to \(K_0\) by real exponentials. Call \(\phi_g\) the resulting right potentials of zero mean (also containing \(m V_R-K_{RL'}n_{L'}\), nonzero modes). They are analytic out to the indicated inverse contours. All zero terms are kept separate.

The current calculation uses the classical level-one lattice realization of the affine algebra~\cite{FrenkelKac1980,Segal1981,Frenkel1985}. We give the mode calculation in our normalization, including the action of the Weyl element on momentum vacua. These signs determine the reflected interaction and its scalar normalization.

Here are details of the current manipulation implementing this Gaussian integral. Use the level-one boson with modes \([H_n,H_m]=2n\delta_{n,-m}\), creation operators \(H_{-n}\), \(n>0\), momentum \(H_0=h\in{\mathbb Z}\), oscillator degree \(N\), and orthonormal momentum shifts \(\sigma_q|h\rangle=|h+2q\rangle\), commuting with oscillators. Concretely the oscillator space is polynomials in commuting \(H_{-n}\) on a unit vector, \(H_n=2n\,\partial_{H_{-n}}=H_{-n}^*\). Write \(\zeta=e^{ic_* s}\),
\[
 E_q(\zeta)=\exp(q\sum_{\nu>0}H_{-\nu}\zeta^\nu/\nu)
             \exp(-q\sum_{\nu>0}H_\nu\zeta^{-\nu}/\nu)
             \sigma_q \zeta^{qH_0+1},\qquad H(\zeta)=\sum H_n \zeta^{-n}.
\]
Notation \(E_q[f],H[f]\) means integration against \(f(s)ds/T\). The modes satisfy \([H[f],E_q(\zeta)]=2q f(\zeta)E_q(\zeta)\),
\([E_{+,n},E_{-,m}]=H_{n+m}+n\delta_{n,-m}\), \(E_{q,n}^*=E_{-q,-n}\). Indeed multiplication and normal ordering gives the factor \((1-\eta/\zeta)^{-2}\) in \(E_+(\zeta)E_-(\eta)\); the pole terms are \(\zeta\eta/(\zeta-\eta)^2+\eta H(\eta)/(\zeta-\eta)\); reversing the mode expansion takes the residue. Same-sign products have no commutator. Thus constant modes act by ordinary \(\mathfrak{sl}_2\) on each finite-dimensional space of fixed \(N+H_0^2/4\). Their Weyl rotation \(W_q=\exp E_q[1]\exp(-E_{-q}[1])\exp E_q[1]\) flips \(H_n\) and sends \(|h\rangle\) to \((-1)^{\lfloor(1+q h)/2\rfloor}|-h\rangle\). For the sign, it is \(+\) on \(|0\rangle,|-q\rangle\), and alternates in steps of 2 by conjugating \(E_q\) to \(-E_{-q}\). The ordinary \(2\times2\) identity (or the finite-dimensional \(\mathfrak{sl}_2\) action) conjugated by \(\exp(qH[\phi]/2)\) gives
\begin{equation}
 \exp E_q[e^\phi]=\exp E_{-q}[e^{-\phi}]\,\exp(qH[\phi]) W_q\,\exp E_{-q}[e^{-\phi}].
 \label{eq:legacy-P-P1}
\end{equation}

Take the normalized \emph{oscillator} trace of the two insertions \(\exp E_{q_g}[e^{\phi_g}]\), \(+\) to left of \(-\), separated on both sides by \(\exp(-bc_* N)\). At each insertion take separately its zero-mode matrix element with
\[
 h_{\rm in}=H_g-q_g k_g,\qquad h_{\rm out}=H_g+q_g k_g,\qquad
 H_g=l-q_g P_g,\quad P_g=\sum_{L'}W_{g j}n_j,\quad k_g=K_{\rm count}-L_g .
\]
(Here \(K_{\rm count}\) denotes the earlier integer \(K\), distinct from the symbol.) Wick contraction produces exactly \(k_g\) currents of type \(g\), the nonzero kernel \(K_0\), signs \((-1)^{\binom{k_g}2}\), and phase coefficients \(q_g H_g\) multiplying \(i c_*s\). The phases follow by replacing \((1-\eta/\zeta)^2\) between currents on the same line by the positive square times \(-\eta/\zeta\). Thermal oscillator contraction uses \(\langle H_n H_{-n}\rangle=2n/(1-e^{-2bc_*n})\) (\(n>0\)); the exponential contraction formula follows by summing the diagonal power series on each one-variable oscillator, or moving positive modes through the trace by geometric series. The right normalizations are exact: if \({\cal L}^*\) denotes the regular diagonal of \(T^{-1}\sum_{p\ne0}K_{gg}(p)e^{-ips}\), then
\[
 2e^{-T/6} e^{k_{gg}/(2T)}\exp({\cal L}^*/2)\tau_T=1
\]
by the collision limit defining \(\tau_T\). Thus one uses the \textbf{full} quadratic energy including regularized diagonal and measure \(ds/(2\pi)\). For the oscillator kernel alone that same convention gives precisely its single-current thermal contraction with measure \(ds/T\) (the vacuum collision slope from \(1-\eta/\zeta\) is \(2\pi/T\)). These rules apply also to unfused left and inverse charges. At a fusion simple pole the self constants add with the regular part of the mutual interaction (whose exponential singularity has unit residue before the regular part is included); the diagonal regularization for the combined species thus gives exactly the local residue up to a unit phase.

In detail, \(\sum_{i<j}q_iq_j\) is minus one half the elementary count. On adding regular self terms to both the nonzero kernel and \(k_{ij}\), the remaining one-body normalizations therefore cancel by the displayed collision limit. For \(K_0\), the regular diagonal is \(-2\log(2\pi/T)+4\sum_{\nu>0} h/(\nu(1-h))\), \(h=e^{-2bc_*\nu}\), yielding the normalized measure and thermal self directly.

Now apply \eqref{eq:legacy-P-P1} and move the left/right inverse currents to \(s+ia_-,s+ia_+\), respectively. Passing the two reflections through the oscillators changes inverse charges to \(t=-,+\) on both lines and Cartan sources to \(\phi_g\) with positive sign. Wick then gives:

\begin{itemize}

\item 

Cartan quadratic covariance \(K_0^{-1}\);

\item 

net linear evaluation on inverse current \(d K_0^{-1}\phi\), including the factor \(e^{-\phi_g(s+ia_i)}\);

\item 

inverse kernel \(B-d K_0^{-1}d(-p)^t\).
These are identities harmonic by harmonic with the preceding trace contractions. For example the same-line positive-\(p\) Cartan-inverse multiplier on including \(-\phi_g\) is \(-t e^{a_i p}\coth(bp)\), the other line gives \(-t e^{a_i p}\operatorname{csch}(bp)\); in the inverse-inverse multiplier the ordering of unequal offsets on one line subtracts \(\operatorname{sgn}(a_i-a_j)\) from \(\coth(bp)\). Gaussian integration replaces \(K_0^{-1}\) by \(K^{-1}\) in these formulas, since the covariance of the field added to \(\phi\) is \(K_0-K\). This gives the nonzero-mode Schur energy, potential, and determinant in \eqref{eq:legacy-P-star}.

\end{itemize}

Indeed for \(p=c_*\nu>0\) the Cartan contraction of \(H_\nu,H_{-\nu}\), within a single Cartan exponential (thus symmetrized) or to the other exponential, is respectively \(\nu\coth(bp),\nu\,{\rm csch}(bp)\). Against an earlier inverse insertion the same-line unsigned factor instead uses \(2\nu/(1-e^{-2bp})\) for the positive mode on the insertion, and against a later insertion this has the extra factor \(e^{-2bp}\); negative modes have reversed rules. This gives the stated multipliers by the mode coefficients of \(E_t\).

\paragraph{Charge matching and zero modes.}
Zero-mode matching matters here. Conservation at \eqref{eq:legacy-P-P1} requires \(\sum_{i:g_i=g}t_i n_i=q_g H_g\), the stated winding constraint. The Weyl sign contributes exponent \(\lfloor(1-k_g-\sum_{i:g_i=g}n_i)/2\rfloor\) to \((-1)\); together with \(k_gP_g+\binom{k_g}2\) this gives a total sign independent of the common integer \(K_{\rm count}\) by the constraint. Inverse phases after same-side symmetrization at group \(g\) have coefficients \(t(k_g+n_{g,-t})\) multiplying \(i c_*(s+i a_i)\). All zero terms linear or quadratic in \(m\) now combine into the displayed finite-part scalar and chemical potentials, on writing \(v=(k_g)_g-m{\bf1}\).

For clarity the rest of the zero energy (before \(r\)) is exactly
\[
 \frac{-1}{2T}\operatorname{FP}_{p=0}\gamma(-p)^t{\cal M}(p)\gamma(p),\qquad
 \gamma(p)=(v+p v',n(p)),\qquad q^t\gamma'(0)=A_L+\pi l/2
\]
where \(q=(q_R,q_{L'},0)\) (counts summed over a double), \(A_L=\sum_{L'}q_j a_j n_j\) and \(v'\) can be complex. Indeed neutrality gives \(q^t\gamma(0)=0\); the two leading pole matrices of \({\cal M}\) are \(4qq^t\) and \(4(qa\,q^t-q(qa)^t)-2\pi{\cal W}\), with \({\cal W}\) given by \(W\) on \(R,L'\), \({\cal W}_{R,i}={\bf e}_{g_i}t_i\), \({\cal W}_{ij}={\bf1}_{g_i=g_j}\operatorname{sgn}(a_i-a_j)t_i t_j\) for inverses. Writing \(\beta=q^t\gamma'(0)\), this contributes (in units \(1/T\))
\[
 2\beta^2-4A_L\beta-2\pi l\beta
 +2\pi i l\sum_{L'}q_j n_j s_j-2\pi i\sum_{L',I'} n_j s_j({\cal W}\gamma)_j .
\]
We write sums schematically treating each individual position. The first
three terms are the shifted center Gaussian \(-(2A_L+\pi l)^2/2\);
the others are exactly the pair/center phases including inverse phases.
The constant matrix supplies \(k_{ij}\) and the corresponding inverse
offsets. The Laurent saddle \(v_*(p)=-QJn(p)\) is regular by the winding
constraint. Matching its leading powers gives the same two conditions,
\[
 q^t(v_*(0),n)=0,\qquad
 q^t(v_*'(0),n'(0))=A_L+\pi l/2.
\]
Thus completion of the square leaves the finite-part Schur term and
\(-h_*(\bar v-u_*)^2/T\). Since
\(\bar v=K_{\rm count}-m-\bar L\), the remaining sum is explicitly
\begin{equation}\label{eq:gas-zero-poisson}
\begin{split}
 &\sum_{K_{\rm count}\in\mathbb Z}
 e^{-h_*(\bar v-u_*)^2/T-r(\bar v+\bar L)}\\
 &\hspace{1em}=\sqrt{\frac{\pi T}{h_*}}
 \sum_{j\in\mathbb Z}
 e^{T(r+2\pi i j)^2/(4h_*)-(r+2\pi i j)(u_*+\bar L)}.
\end{split}
\end{equation}
The integer shift by \(m\) contributes only
\(e^{-2\pi i j m}=1\). This gives both the second square root in
\eqref{eq:gas-transform-scalar} and the last two terms in
\eqref{eq:legacy-P-star}; extracting a fixed Fourier coefficient reverses
this Poisson summation.

More explicitly first use \(\gamma_{\rm raw}=\gamma+(m{\bf1},0)\) (right value \(k\)). Include the source \(m\,\gamma_{\rm raw}(-p)^t\allowbreak(V_R,V_{L'},0)\) inside the Laurent prescription, in \(1/T\) units. Its pole gives the phase slopes \(-2\pi i m W^V_i s_i\) on the left (\(W^V_R=0\)); its constants supply the potential means. The preceding quadratic calculation applies unchanged to \(\gamma_{\rm raw}\); in particular \(({\cal W}\gamma_{\rm raw})_i=-t_i(k_g+n_{g,-t_i})\) on inverses, and the constant right-inverse and inverse-inverse entries at fixed group are \(2\pi t_i a_i,-2\pi|a_i-a_j|t_i t_j\), giving the rest of their offset phases. Completing the square with \(V_R=K{\bf1}\widehat\rho\) now separates the scalar and chemical zero modes as asserted, leaving \(\gamma\) here (\(m\widehat\rho\,{\bf1}\) can be subtracted instead of \(m{\bf1}\), with the same finite part since the rows against \({\bf1}\) in the right block have at worst a simple pole). For the saddle, the common projection of \(Jn(p)\) has no pole by the constraint, while the difference channel of \(K\) has eigenvalue \(8/p^2+O(1)\); matching leading two orders gives \(q^t(v_*(0),n)=0,\ q^t(v_*'(0),n'(0))=A_L+\pi l/2\). The sign independence used here follows also by incrementing \(k_g\) at both insertions: the change at each is \(P_g+\sum_{i:g_i=g}n_i=l q_g\bmod 2\).

\subsection{Convergence of the current identity}
The preceding mode computation must hold for actual analytic tests, not only formal series. We establish absolute convergence before moving the remaining contours.

Here are further details justifying the trace argument. Work first with fixed integer coefficient (fixed \(K_{\rm count}\)), fixed \(l,L'\) counts and left positions, and real \(T\). We can also do this for \(k_g<0\), taking the original right coefficient as zero via its zero-mode matrix element. The right mean-zero Gaussian field has almost surely analytic samples on closed strips of width less than \(L/2\) on either side, by summability of the square roots of its Fourier variances with exponential weights. The deterministic terms of \(\phi\) are also analytic on strips reaching strictly past the inverse contours: in the displayed shifted factors coupling to the preliminary left contours each noncancelled zero or pole has distance at least \(9L/20\), and the potential is analytic there. One can therefore first do the trace identity for trigonometric polynomial \(\phi\), inserting a formal scalar multiple \(\lambda\phi\).

Indeed coefficients at fixed order in \(\lambda\) connect only bounded ranges of differences in the level \(N+h^2/4\) (since \([N+h^2/4,E_{q,n}]=-n E_{q,n}\) for the coefficient of \(\zeta^{-n}\)). The constant modes are mutually adjoint raising/lowering operators of norm \(O(\sqrt{1+N+h^2})\) on each level by the finite-dimensional unitary spin representation (the weights of \(H_0\) there are bounded by twice the square root of the level). Fixed modes likewise have polynomially bounded norms: \(E_{+,n}, E_{-,-n}\) generate the same adjoint relations with weight \(H_0+n\) on finite-dimensional eigenspaces of \(N+h^2/4+nH_0/2\), and the oscillator modes act by differentiation and multiplication. These bounds follow also from the ordinary raising/lowering norm recursion on each highest-weight string, \(\|E f\|^2-\|F f\|^2=-\langle f,H f\rangle\) for weight commutator \([E,F]=H\), \(F=E^*\). Thus the coefficient identities obtained by conjugation as above can be used under the projected thermal trace, including termwise expansion of constant-mode exponentials (norm sums on energy \(d\) bounded by \(C(1+d)^C\exp(C\sqrt{d+1})\) at each fixed Taylor order, by summing the unperturbed runs of constant modes in the exponential series; momentum projections between the two rays reset only fixed endpoint momenta). The oscillator multiplicities are subexponential by the convergent generating function at any positive damping. At each fixed number of currents the Wick calculation gives precisely these mode integrals: use first slightly separated radii in cyclic radial order, then identify same-side radii (zeros not poles). Termwise traces of normal-ordered oscillator series at strictly separated radii within the thermal period converge absolutely, for finitely many insertions and Cartan coefficients: replacing all oscillator coefficients by absolute values on the occupation basis still gives summable geometric-series contractions by strict separation. Radius displacements preserve the integrated mode products. This proves the claimed trace transform coefficientwise in \(\lambda\).

We detail the convergence bound allowing actual evaluation of this identity. In summing the inverse integrals with fixed other data, even averaging termwise \emph{absolute} values over the Gaussian after Wick is allowed. The real symmetric nonzero-mode inverse block controlling this sum (Cartan real quadratic HS tilt completed by the Gaussian) is, omitting the positive factor \(2\pi/p\), in each color eigenspace for \(p>0\),
\[
 \begin{gathered}
 t t'\left[\Lambda\cosh((a_t-a_{t'})p)-\sinh(|a_t-a_{t'}|p)
 -\frac e{1-e}\Lambda\sinh(a_t p)\sinh(a_{t'}p)\right],\\
 \Lambda\in\{\coth(bp/2),\tanh(bp/2)\}.
 \end{gathered}
\]
Indeed only the even part of \(d\), with rows \(-(2\pi/p)t\sinh(a_t p){\bf e}_{g}^t\), couples in real part (times \(K_0^{-1}\)) to the real field, whose covariance after quadratic tilt is \((e/(1-e))K_0\). The first two terms are the symmetrized \(B-dK_0^{-1}d(-p)^t\). Deterministic linear sources cost at most linear growth in inverse counts plus constants (fixed other data), by Fourier decay. For a deterministic \(\lambda\phi\) without Gaussian integration simply omit the last block term, uniformly on compact \(\lambda\) sets.

This matrix is positive definite. Write \(x=Lp,\ A=5x/12,B'=-x/4,U=A-B'=2x/3\). We have \(e/(1-e)\le e^{-x}\), equivalently \((1+e^x)(\sinh3x-\sinh2x)\le\sinh4x\). Diagonals are then positive. The determinant divided by \(\Lambda\sinh U\) is
\[
 2\sinh(3x-U)/\sinh(3x)-
 (e/(1-e))\,[\Lambda\sinh U-\cosh U+\cosh(A+B')].
\]
Here \(\Lambda\sinh U<\cosh U\), and \(2\sinh(3x-U)/\sinh(3x)\ge(14/9)e^{-2x/3}\) dominates the required term. At high frequency the matrix tends to identity with exponential error. Thus the regularized full nonzero-mode real energy contributes at most \(O_T(M\log(M+2))\) to log absolute weights for \(M\) inverse particles. Indeed smooth with a factor \(e^{-|p|/(1+M)}\) and keep also the diagonals in the nonnegative quadratic form. Undoing the smoothing costs at most \(O_T(1/(1+M))\) per off-diagonal pair in log upper bound: the decaying remainder is smooth and the only singular log interactions are repulsive \(2\log|1-e^{ic_* s}|\), using \(|1-e^{i\theta}|\le e^{\epsilon/2}|1-e^{-\epsilon+i\theta}|\). The regularized self discrepancy is logarithmic. On the other hand the inverse offset zero-mode phases give
\[
 -c_*(a_+-a_-)\sum_g n_{g,+}n_{g,-}+O(M+1)
\]
with fixed other data and fixed \(n_{g,+}-n_{g,-}\). This is strictly quadratically negative. This proves absolute convergence, entireness in the deterministic test multiplier used above, and also convergence under trigonometric approximation of analytic \(\phi\) on the indicated strips (before HS averaging). Wick Gaussian completions of the convergent modes follow by finite cutoffs and uniform Gaussian integrability (the tilt eigenvalues \(e(p)\) are strictly below 1 and summable); the same absolute estimate gives Fubini for the averaged sum. For fixed counts, coefficient and \(l\), the deterministic bounds here on the initial lines can be taken uniformly in the left positions (bound the right-left couplings on a closed analytic strip containing the needed offsets; the purely left weights themselves have no uncancelled poles on the integration torus after initial double bindings).

\subsection{The final contours and their coercive kernel}
\paragraph{Contour continuation and fusion.}
The next contour move creates one additional bound species. We give its exact symbol before proving the real quadratic estimate that controls the grand sum.

For estimates we continue the \(D\) line upward, collecting additional residues. Here are explicit final mode rules for \eqref{eq:legacy-P-star}; \(p\) remains in physical length units. The species in each color are indexed \(i=0,1,2,3,4\) (respectively \(S,D,I'_+,I'_-,F\), with \(F\) a bound same-color \(D,I'_-\) pair). Set
\[
 x=Lp,\quad z=e^x,\quad
 \delta=(6,1,5,-3,-1)/12,\quad w=(1,1,-1,1,2)^t,\qquad b_0=(2),\ b_1=(2,3).
\]
Thus \(\delta L\) are the imaginary offsets in the center coordinates, whose physical elementary offsets are obtained by adding \(b_i L\) on \(S,D\) (and the constituent \(b_1L\) on the \(D\) member of \(F\)). In formulas for color matrices we give the scalar rules for the two color eigenspaces \((1,c)\), \(c=1,-1\). Define
\[
 H(a;\sigma)=\sum_{j=0}^7 h_j\frac{z^{[a+j]_8}}{z^8-1}+\tfrac{1-\sigma}{2}h_{[-a]_8},
 \quad h=(2,-1,1,-2c,2c,-2c,1,-1)
\]
at integer \(a\), using standard residues \(0,\ldots,7\) for the brackets. Write \(H_*=H(0;0)\), and set
\[
 A_0=\sum_{a\in b_0}H(a;1),\quad A_1=\sum_{a\in b_1}H(a;-1),
 \quad A_2=1,\quad A_3=-1,\quad A_4=A_1+A_3 .
\]
For \(0\le i,j<4\) let \(P_{ij}\) be
\[
 \begin{cases}\displaystyle\sum_{a\in b_i,k\in b_j}H(a-k;\operatorname{sgn}(j-i)),& i,j<2,\\
 \operatorname{sgn}(j-i),&i,j\ge2,\\0,&\text{otherwise}.\end{cases}
\]
Then put
\[
 \begin{split}
 F_{ij}^{(c)}(z)&=P_{ij}(z)+A_i(z)A_j(z^{-1})/H_*(z)
 +(-1+2c)\operatorname{sgn}(i-j){\bf1}_{\{i,j\}=\{1,3\}},\qquad i,j<4,\\
 F_{i4}&=F_{i1}+F_{i3}+{\bf1}_{i=3}-{\bf1}_{i=1},\qquad
 F_{4i}(z)=-F_{i4}(z^{-1}),\qquad F_{44}=F_{11}+F_{13}+F_{31}+F_{33}.
 \end{split}
\]
The final kernel is
\begin{equation}\label{eq:gas-final-kernel}
 Z^{(c)}_{ij}(p)=(2\pi/p)e^{(\delta_i-\delta_j)x}F^{(c)}_{ij}(z).
\end{equation}
Other final rules (\(A_i\) here at \(c=1\)) are
\begin{equation}\label{eq:gas-final-rules}
 \begin{split}
 &\sum w_i=0,\qquad
 u_*=\frac{\pi}{h_*}\lim_{p\to0}\sum \frac1p A_i(z^{-1})e^{-\delta_i x+i p s},\qquad
 2\bar L=\sum (1,2,0,0,2)_i,\\
 &\mu_i(p)=-\frac{2\pi}{p} (0,1,1,-1,0)_i e^{\delta_i x}\widehat\rho(p).
 \end{split}
\end{equation}
The sums in the first line are over particles. In particular \(A_i(1)=-w_i\). The chemical rule follows by substituting the potential array in the mode formulas before fusion and \(\widehat\rho=1/(z+z^{-1}-1)\).

Here is the continuation check. Initially omit species 4, take \(\delta_1=-9/20\) only changed entry, and omit the correction supported on \(\{1,3\}\). The formulas then give exactly the Schur kernel and other terms above, by extracting the offset exponentials in \(\mathcal M\). All the rational functions after offset extraction in the kernel are bounded at \(z=0,\infty\), with skew reflection \(F_{ji}(z)=-F_{ij}(1/z)\). For instance at the \emph{final} branches on \(i,j<4\) their matrix at \(\infty\) by the displayed formulas is
\[
 \begin{pmatrix}
 1&0&0&0\\2c&1&1-2c&0\\-1&0&1&0\\1&-1+2c&-2&1
 \end{pmatrix}.
\]
To see the initial extraction, the row against the right block after pulling out \((2\pi/p)e^{\delta_i x}\) is precisely \(A_i\): a shift at an integer multiple of 8 in the sum defining \(H\) approached with a negative offset has exponent index 8 rather than 0, giving the extra term for sign \(-1\). The column then uses \(-A_i(z^{-1})\), and the same shift prescription between left constituents gives \(P\). In the common projection only, \(A_i(1)\) is also the column sum of the right winding entries, so their integrality (from the shifted arrays with no exact boundaries against \(R\)) gives equivalence of \(\sum w_i=0\) to the constraint including integral \(l\). For the chemical term replacing \(h\) in the two left sums defining \(A_i\) by the potential array gives \(((0,1)_i-A_i)\widehat\rho\) at \(c=1\), by substitution.

Thus after removing the high-frequency constants the pair series are analytic throughout the contour moves (offset differences strictly between \(-1,1\)). For offsets initially strictly ordered, say \(\delta_i<\delta_j\), the high-frequency logarithm with coefficient \(\nu\) in the unsplit color basis gives a factor
\((1-e^{-ic_*d})^\nu\), where \(d=s-s'+iL(\delta_i-\delta_j)\). In carrying the lower contour to above the other, subtract \(\nu\) from the ordered coefficient function \(F_{ij}\), adding to \(F_{ji}\), times the given color matrix unit(s). The resulting expansion uses \((1-e^{ic_*d})^\nu\) instead. Their ratio is compensated up to a unit constant phase by the change at Laurent finite part zero (from the coefficient \(2\pi\nu/p\) times \(e^{-ipd}\)). Equal-offset logarithms of even order can likewise use the averaged branch, with the regularized diagonal as before. All such phase choices here are independent of \(K_{\rm count}\).

In moving \(D\) up the crossed interaction \(D,I'_-\) thus has order \(-1\) in the same color, 2 across colors. Other crossed interactions from moving multiple lines of a given species relative to each other do not give poles. Each resulting simple-pole residue is \(F\). Its interactions have positive orders with \(D,I'_-\) (summing the constituent orders including the same-color same-species squares) and between residues, so only disjoint matchings contribute and the bound species itself moves freely to the indicated contour. The correction on \(\{1,3\}\), and the formulas with 4 (summing the constituent entries, adjusting for the \(F\) offset intermediate between \(D,I'_-\) for formerly equal-offset pairs), give precisely the resulting branch prescription. Its self term and measure follow by regularized collision as for the earlier fusion, and the matchings leave the factorial measures in final species. These deformations can all be made termwise with fixed integer coefficient before Poisson summation in \(K_{\rm count}\): each integrand then is periodic before taking residues, as also clear from the integer current/winding phases before compressing the Laurent zero term. The adjustments concern only the Schur log pair functions, not the transported \(u_*\), and bring no integer-coefficient-dependent phase. Consequently, once absolute convergence of the final sum \eqref{eq:gas-grand-sum} on a full imaginary \(r\)-period has been shown, its identity follows by integrating against each Fourier monomial. Indeed this recovers the fixed-coefficient Gaussian with its unchanged phases. Group by the original left counts (a bound particle uses one \(D\)) and \(l\); for each such group the preliminary inverse expansion was absolutely convergent and contour deformation agrees term by term, and groups with negative required right counts integrate to zero. Remaining groups recover precisely the opened original coefficients.

In detail the poles being bound are at equality of the \emph{shifted center} coordinates (mod period); the one-body symbols have no singular crossings here. The orders with the exterior \(1,3\) are then 1 (same color) and 2 (other), and between two bound particles are 2 and 4. Thus contours can for example be moved successively. For \(F_{14}\) the unadjusted positive-frequency constant of the constituent sum is 1 which must now go to zero (row above column); for \(F_{34}\) it is the corresponding \emph{transposed} constant which must go from 1 to zero, giving the other adjustment. No adjustment of the constituent sum in \(F_{44}\) is needed. Its regular diagonal in same color adds both old regular diagonals and twice the pole regular part (up to phase if evaluating from one side); indeed the logarithmic singular terms add with orders \(2+2-2=2\), and the regular constants from opposite determinations of the simple pole differ at most by a fixed imaginary phase, as seen from \(1-e^{\pm i c_*d}\) at 0. The residue with the relative-coordinate measure \(dd/(2\pi)\) therefore realizes this prescription, with at most a fixed unit phase. Branch changes on exterior factors likewise have only the fixed phase after finite-part compensation (the exact ratio of the two geometric arguments is \(-e^{-ic_*d}\)). Thus all phase choices here can be made independently of \(T\) as well as of \(K_{\rm count}\).

There is no additional phase depending on the real ordering of the
particles. Different final species have distinct offsets, so exchanging
their real positions crosses no contact singularity. At a common offset,
the nonzero contact orders are 2 for the same species and color, and 4
for \(F\) in opposite colors; these are even. Reversing the real order
therefore introduces no further branch sign. The remaining contour,
Weyl, and fusion constants depend only on the multiplicities (and
possibly on \(m\)).

We can now specify the complete representation whose absolute value will
be estimated. Its phase factors are the ones just transported through
the contour moves.

\begin{auxproposition}[Current representation on the final contours]
For every integer \(m\ge1\) and all sufficiently large real \(T\),
\(P_{m,T}(e^r)\) equals \eqref{eq:gas-grand-sum}, with
\(\alpha=(g,i)\) ranging over the five final species in both colors and
\[
 \nu_{g,i}\in\mathbb Z_{\ge0},\qquad
 \sum_g\sum_{i=0}^4w_i\nu_{g,i}=0.
\]
The exponent uses \eqref{eq:gas-final-kernel} and
\eqref{eq:gas-final-rules}, with the Laurent and regular-diagonal
conventions preceding \eqref{eq:gas-transform-scalar}. The scalar is
\eqref{eq:gas-transform-scalar}; the factorial measures and the sum over
\(j\in\mathbb Z\) are exactly those in \eqref{eq:gas-grand-sum}.
For the fixed branches in the derivation,
\(|\varepsilon_m(\nu)|=1\), and these factors depend only on \(m\)
and the multiplicities, not on the positions, \(T\), \(r\), or \(j\).
The sum and integrals converge absolutely, locally uniformly in \(r\).
\end{auxproposition}

The coefficient identity and phase dependence have been established
above. To complete this proposition, we next prove the energy estimate
that permits summing all multiplicities and all \(j\). We will then
continue the representation in \(T\) and localize the particles near
the period ends.

\paragraph{The real kernel bound.}
Put \(\eta=.55,\ \alpha=2\pi^2/h_*\). In full species-color space
(index color first), write \(W=(w,w)^t\) as one column, and
\(Z_{\rm s}=(Z+Z^t)/2\).

\begin{auxlemma}[Coercivity and frequency behavior of the final kernel]
There is \(\epsilon_0>0\) such that, for every real \(p\ne0\),
\begin{equation}
 Z_{\rm s}(p)\ \succeq\ \eta\alpha\, WW^t/p^2 + P_{\times}(p) + \epsilon_0(1+p^2)^{-1}{\bf I}
 \label{eq:legacy-P-P2}
\end{equation}
where \(P_{\times}\) has just equal off-diagonal color blocks given by
\[
 \frac{2\pi}{p}\left(\frac{11}{12}\frac{2\sinh(x/2)}{2\cosh x-1}hh^t
       +\tfrac32\tanh(x/4) f f^t\right),\quad
 h=(4,5,0,4,9)^t/5,\quad f={\bf e}_4 .
\]
Moreover,
\begin{equation}\label{eq:gas-kernel-zero}
 Z(p)=\alpha WW^t/p^2+M_*/p+O(1)\qquad(p\to0),
\end{equation}
where \(M_*\) is real and skew. At high real frequency, \(Z(p)\) is
\(2\pi/|p|\) times the identity plus the off-diagonal color coupling
\(2ff^t\) in the same units, up to exponentially decaying error.
\end{auxlemma}

We prove the coercive bound by reducing it to two finite polynomial
matrix inequalities; the frequency assertions follow directly from the
same rational symbols. Use \(p>0,\ t=e^{x/12}\) (the symmetric kernels
in \eqref{eq:legacy-P-P2} are even). Form the symmetric \(5\times5\)
matrix \(Y_c\) from
\(t^{12(\delta_i-\delta_j)}F^{(c)}_{ij}(t^{12})\) by averaging with
its transpose, then subtract
\[
 c\,[p_1hh^t+p_2ff^t]+{\bf1}_{c=1}p_3ww^t,
 \quad p_1=\tfrac{11}{12}\frac{t^6(t^{12}-1)}{t^{24}-t^{12}+1},
 \quad p_2=\tfrac32\frac{t^6-1}{t^6+1},\quad p_3=\frac{t^6+1}{8(t^6-1)} .
\]
\paragraph{The matrix positivity certificate.}
Use the positive multipliers given by
\[
 \begin{gathered}
 d_+(z)=(z-1)(z^2-z+1)(z^2+z+1)^2,\qquad
 d_-(z)=(z+1)(z^2-z+1)^3,\\
 N=(1+t^6)^{(1-c)/2}d_c(t^{12}),\qquad M=\deg(N)/2
 \end{gathered}
\]
(so \(M=42,45\) respectively). Then \(N Y_c\) is polynomial with its matrices of coefficients at indices \(M\pm j\) related by parity \(c\). Write \(B_j=[t^{M+j}]N Y_c\), \(0\le j\le M\), with the sole modification of \textbf{halving} \(B_0\) when \(c=+1\). Put
\[
 v=\tfrac{M}{6}\log t,\quad a=2\cosh(v)-2,\quad \beta_j=6j/M .
\]
Thus \(Nt^{-M}Y_+\) or \(Nt^{-M}Y_-/(2\sinh v)\), respectively, is \(\sum_j B_j\) times the respective weights (term by term)
\[
 2\cosh(\beta_j v),\qquad \sinh(\beta_j v)/\sinh v .
\]

Use the following short polynomial comparison. For \(c=+1\) take \(K=6\), and for \(c=-1\) take \(K=5\). Define coefficients \(g_k\) at \(\beta\) by
\[
 g_k^+(\beta)=\frac{2}{(2k)!}\prod_{r=0}^{k-1}(\beta^2-r^2),\qquad
 g_k^-(\beta)=\frac{\beta}{(2k+1)!}\prod_{r=1}^{k}(\beta^2-r^2).
\]
For integer \(0\le\beta\le6\) each weight is exactly \(\sum_{k=0}^K g_k(\beta)a^k\). For instance the differential equations in \(a\) for the two functions are \(a(4+a)y''+(2+a)y'=\beta^2 y\) and \(a(4+a)y''+3(2+a)y'=(\beta^2-1)y\); Taylor recurrence gives the formulas, terminating at integers. For other \(0<\beta<6\) the true value lies between the two truncations at \(K-1\) and \(K\) for every \(v>0\). Indeed at fixed \(v\) these are interpolation in \(\beta^2\) at the squares of successive integers starting from 0, or for minus at squares starting from 1 after division by \(\beta\). All derivatives of the interpolated function are positive by its series. Hence the remainder has sign of the node product by repeated Rolle, and the last node \(6^2\) changes its sign.

This proves the lower bound \(\sum_{k=0}^K U_k a^k\) (in positive semidefinite order) for the respective cleared matrix, with
\[
 \begin{split}
 U_k&=\sum_{j=0}^M g_k(\beta_j) B_j\qquad(k<K),\\
 U_K&=\sum_{j=0}^M g_K(\beta_j) B_j-
 \tfrac12\sum_{\beta_j\notin\mathbb Z}
 \left[g_K(\beta_j)B_j+|g_K(\beta_j)|\operatorname{diag}_i\Big(\sum_n |(B_j)_{in}|\Big)\right].
 \end{split}
\]
This uses the ordinary row-sum diagonal bound to either side of a symmetric matrix, allowing half the last term as midpoint when nonintegral.

For clarity the small coefficient inequalities used next are given explicitly below, with details for verifying them. In the plus case \(U_0=(7/2)ww^t\). For the other coefficients the following symmetric matrices \(J_k\) approximate \(100 U_k\) with elementwise error \(<1\) (successive upper-triangle rows separated here by semicolons).

{\scriptsize
\begin{verbatim}
c=+1
k=1: 1745 410 -2105 1011 1062; 1460 -783 -863 238; 3885 -2644 -3508; 2333 1111; 1402
k=2: 6082 977 -4261 2809 2141; 4542 -882 -2052 941; 9684 -4638 -5401; 6393 2792; 4079
k=3: 7166 829 -3563 2893 1521; 5808 -433 -1507 2090; 9655 -3094 -3107; 7193 3475; 6340
k=4: 3787 176 -1414 1213 197; 3355 -48 -554 1401; 4554 -784 -520; 3786 1833; 3948
k=5: 927 -26 -265 202 -92; 877 27 -106 358; 1016 -37 76; 927 408; 1046
k=6: 50 -11 -3 -3 -23; 45 10 -11 1; 63 10 24; 37 1; 32
\end{verbatim}
}

{\scriptsize
\begin{verbatim}
c=-1
k=0: 188 48 27 -52 102; 933 747 -619 208; 640 -427 160; 828 102; 430
k=1: 1435 187 -198 -522 556; 5710 4260 -4101 954; 4109 -2744 464; 5134 377; 2433
k=2: 3840 126 -1127 -1441 902; 9533 6247 -6689 1216; 6906 -4144 192; 8587 271; 4310
k=3: 4041 -132 -1354 -1314 397; 5635 3085 -3271 284; 4430 -1835 14; 5201 -150; 2516
k=4: 1229 -63 -421 -360 50; 1331 544 -538 -61; 1161 -267 -1; 1269 -122; 594
k=5: 55 -2 -17 -12 0; 53 10 -9 -8; 66 -3 0; 53 -8; 15
\end{verbatim}
}

Here is one arithmetic recipe for the table. The upper entries of \(D_{ij}=d_c F_{ij}^{(c)}\) for \(i,j<4\) by direct substitution are polynomials with these coefficient strings (ascending, triangular rows):

{\scriptsize
\begin{verbatim}
c=+1
1 1 0 2 2 0 1 1; 2 0 2 1 2 0 1 0; -1 0 -2 -1 -2 0 -2 0; 1 0 2 1 2 0 2 0
1 1 1 1 1 1 1 1; 0 0 -1 -2 -2 -1 -1 -1; 1 0 2 1 3 0 1 0
1 1 1 1 1 1 1 1; -2 -1 -2 0 -2 0 -1 0
1 1 1 1 1 1 1 1
\end{verbatim}
}

{\scriptsize
\begin{verbatim}
c=-1
-1 5 -12 18 -18 12 -5 1; 2 -4 6 -5 4 -2 1 0; 1 -2 4 -5 6 -4 2 0; -1 2 -4 5 -6 4 -2 0
-1 1 -1 1 -1 1 -1 1; 0 0 1 0 0 3 -3 3; 3 -6 8 -3 -3 6 -3 0
-1 1 -1 1 -1 1 -1 1; 2 -3 4 -2 0 2 -1 0
-1 1 -1 1 -1 1 -1 1
\end{verbatim}
}

Fill the rest by \(D_{ji}(z)=c z^7 D_{ij}(z^{-1})\), and extend to index 4 by the same addition rules as for \(F\) (the added constants multiplied by \(d_c\)). These strings also show the claimed degree bounds after offset multiplication. To reproduce the entry bounds put \(o_{in}=12(\delta_i-\delta_n)\). For each entry take the successive coefficients at powers \(M+j\), \(0\le j\le M\), in
\[
 \frac{(1+t^6)^{(1-c)/2}}{2}
 (t^{o_{in}}D_{in}(t^{12})+t^{-o_{in}}D_{ni}(t^{12}))-
 N[c h_i h_n p_1+c f_i f_n p_2+{\bf1}_{c=1}w_i w_n p_3],
\]
with \(j=0\) halved for plus; products of \(p_r\) in use are polynomial after cancellation with factors of \(N\). Weight this row of numbers at \(j\) by \(100 g_k(6j/M)\) and sum; for \(k=K\) halve the nonintegral-node weights and on diagonal entries also subtract the corresponding absolute row-sums with absolute half weights. These sums rounded to within 1 give the displayed arrays (successive weights \(g_{k+1}/g_k\) when nonzero are \(((6j/M)^2-(k+(1-c)/2)^2)/[(2k+1+(1-c)/2)(2k+2+(1-c)/2)]\)).

Write \(H_k=(J_k-5{\bf I})/100\), except in the plus case take \(H_0=U_0\). Thus always \(U_k\succeq H_k\). The positive-definite checks on the displayed matrices are summarized here as strict lower bounds on \textbf{each} successive leading elimination pivot (starting with entry \(00\)):
\[
\begin{array}{c|l|l}
 c&\text{matrix}&\text{bound}\\\hline
 -&H_0,H_1,\ldots,H_5& 6/100\\
 +&H_2,\ldots,H_6&2/100\\
 +&\displaystyle\sum_{k=0}^i \binom{i}{k}\binom4k^{-1}4^{-k}H_k,\quad i=1,2,3,4&5/100\\
 +&\displaystyle\sum_{k=1}^4 k4^{1-k}H_k&1
\end{array}
\]
These bounds just use the preceding integer arrays with diagonal correction: e.g. starting with the indicated linear sums as \(b_{in}\), take \(b_{dd}\) in order \(d=0,\ldots,4\), replacing the remaining entries by \(b_{in}-b_{id}b_{dn}/b_{dd}\) before the next bound. The matrices with binomial factors are the nonconstant-index Bernstein coefficients of \(\sum_{k=0}^4 H_k a^k\) on \([0,1/4]\) (expand in \(\binom4i(4a)^i(1-4a)^{4-i}\)); the remaining coefficient there is \(H_0\). At \(1/4\) the derivative is positive definite by the last check, increasing thereafter since \(H_2,H_3,H_4\) are positive. This gives strict positivity of the comparison throughout \(a>0\).

Thus \(Y_c\) are strictly positive, with smallest eigenvalue at least constant times \(\tanh x\) (use the linear-order Bernstein bound near zero for plus, \(H_0>0\) there for minus, and the top coefficient near infinity in both). In the plus sector the amount needed inside the dimensionless matrix for the winding term of \eqref{eq:legacy-P-P2} is \((8\eta/9)ww^t/x\), covered by the displayed subtraction since \(\coth(x/4)\ge4/x\). This proves \eqref{eq:legacy-P-P2}.

For \eqref{eq:gas-kernel-zero}, the common-color symbol satisfies
\(H_*=h_*p/(2\pi)+O(p^3)\), the unfactored \(P_{ij}\) terms in
\(F\) are bounded, and the difference-color entries \(F^{(-)}\) are
regular because \(d_-(1)\ne0\). The skew-reflection rule makes the
coefficient \(M_*\) real and skew. At high frequency the displayed
branch constants give the identity and the off-diagonal color coupling
\(2ff^t\); the remaining rational terms decay exponentially. This
completes the kernel bound, including the asymptotics needed to control
contact singularities and to continue the period.

\subsection{Stability, continuation, and the size estimate}
The matrix bound now becomes a spatial energy inequality for every
finite winding-neutral configuration. This estimate has no site-number
parameter or external potential; the interval operators in
Section~\ref{u:interval-operator} will use that independence. We then
add the \(m\)-dependent potential and choose a logarithmic period.

\paragraph{Spatial stability.}
Let \(T=T'\) be any sufficiently large real period, and let
\(s_d\in[-T'/2,T'/2]\) be the positions of \(N_0\) particles of the
final species, with \(W_d=w_{i_d}\) and \(\sum_dW_d=0\).
Define
\[
 \mathcal E_Z=\frac1{T'}\sum_p n(-p)^tZ(p)n(p),
\]
using the Laurent finite part at zero and the regular diagonal at
contact. The same conventions apply to other kernel subscripts.
Distinct-particle coincidences can be omitted as null sets. Put
\[
 {\cal D}(t)=\sum_{s_d>t} W_d,\qquad S_w=\int{\cal D}=\sum W_d s_d,\qquad
 {\cal Q}=\sum_{d,e}\sum_{\nu\in\mathbb Z}(1+|s_d-s_e+\nu T'|^2)^{-1},
\]
where integrals are over the centered period interval.

\begin{auxlemma}[Spatial stability at a real period]
There are absolute constants \(c_1>0\) and \(C_0<\infty\) such that,
for every sufficiently large real \(T'\), every configuration just
defined, and every \(j\in\mathbb Z\),
\begin{equation}
 \Re{\cal E}_Z+\alpha(2jS_w+j^2T')\ \ge\
 \alpha\eta\int({\cal D}+j)^2+c_1{\cal Q}-C_0N_0.
 \label{eq:legacy-P-P3}
\end{equation}
The constants are independent of the period, the particle counts and
positions, and \(j\).
\end{auxlemma}

To prove the lemma, at the zero Laurent mode only shift \(n(p)\)
formally by \(ipjT'v_0\), where \(v_0\) is a real species-color
vector with \(W^tv_0=1\). Write the resulting energy with
superscript \([j]\). Then
\[
 {\cal E}_{Z_{\rm s}}^{[j]}={\cal E}_{Z_{\rm s}}+\alpha(2j S_w+j^2T'),\qquad
 {\cal E}_{WW^t/p^2}^{[j]}=\int({\cal D}+j)^2
\]
by winding neutrality and Fourier series (nonzero coefficients obtained by differentiating \({\cal D}\)). Indeed the finite part of the winding form comes just from its first-order terms in each factor, giving \((S_w+jT')^2/T'\), the constant harmonic contribution on the centered period interval. The real part for the unshifted raw energy is just that of \(Z_{\rm s}\).

The \(P_{\times}\) term has nonnegative energy pointwise: each entry of its off-diagonal color blocks has a nonnegative spatial kernel by periodization. Indeed \({\rm sech}(\lambda p)\) transforms to a positive density proportional to another sech (shift the contour by \(i\pi/\lambda\) and take the one pole), also its square thus has nonnegative transform, and \(\tanh(\lambda p)/p=\int_0^\lambda{\rm sech}^2(up)\,du\). Products transform by convolution. Also \(\sinh(x/2)/x=\int_0^{1/2}\cosh(u x)du\), and \(1/(2\cosh x-1)=\sum_{n\ge1}(2\cosh x)^{-n}\); \(\cosh(ux)/\cosh x\) has positive transform since the two shifted sech transforms are conjugates of positive real part for \(0\le u\le1/2\). This series against \(\cosh(ux)\) converges in integral norm. The resulting periodic kernels with at most logarithmic singularity are evaluated off contact only. The last term in \eqref{eq:legacy-P-P2}, using the inverse kernel \(e^{-|s|}/2\) periodized, has energy at least \(c_0{\cal Q}\), \(c_0>0\). Indeed divide the circle into intervals of lengths between \(1/2\) and 1. The sum of squared interval counts in each species then controls \({\cal Q}\) by summability and Cauchy-Schwarz.

After these subtractions the symbol \(Y_0\) on the left of \eqref{eq:legacy-P-P2} (i.e. the difference of the two sides) is positive semidefinite. It is even, with at most the rank-one \(WW^t\) leading double pole, so neutrality also ensures finite-part positivity (apply positivity at real \(p\to0\), even with the indicated \([j]\) shift). To control the diagonal convention, smooth \(Y_0\) by multiplication by \(m_\lambda(p)=\exp(-\lambda(\sqrt{1+p^2}-1))\), with \(\lambda>0\) small fixed. The smoothed \emph{full} quadratic form is nonnegative. In undoing the smoothing write
\[
 (1-m_\lambda)Y_0
 =2\pi B_0(1-e^{-\lambda|p|})/|p|+E_\lambda(p);
\]
\(B_0\) is identity plus the off-diagonal color coupling \(\tfrac12 ff^t\) (in particular entrywise nonnegative). \(E_\lambda\) has \(L^1\) norm and distributional second derivative total variation both \(o(1)\) as \(\lambda\downarrow0\). Indeed on \(|p|\le1\), \((1-m_\lambda)/p^2\) is analytic smooth of size \(O(\lambda)\) also with two derivatives, and the cusp from \((1-e^{-\lambda|p|})/|p|\) has slope jump \(O(\lambda^2)\). On \(|p|\ge1\) subtract \(2\pi B_0/|p|\) first in \(Y_0\), leaving \(O(1/|p|^2)\) with decaying derivatives by the displayed symbols; multiplying by \(1-m_\lambda\) gives the claimed norm bounds there. The mismatch in cutoff on the leading term is \(O(\lambda e^{-\lambda |p|}/|p|)\), also with two derivatives, since \(\sqrt{1+p^2}-|p|=O(1/|p|)\). Hence periodized \(E_\lambda\) costs \(o(1){\cal Q}\). The other term gives the periodization of \(B_0\log(1+\lambda^2/s^2)\) by elementary Fourier transformation. It is nonnegative off contact; its regularized diagonal is bounded below by \(2\log\lambda\). These periodizations include exact ordinary zero means as required for the discrepancy, which has no pole. Choose \(\lambda\) small enough to absorb the \(o(1){\cal Q}\) term. Its fixed regular-diagonal cost is at most a constant per particle, proving \eqref{eq:legacy-P-P3} uniformly in \(T'\).

This also completes the real-period current representation. For fixed
\(m,T'\), the chemical potential is bounded above by a constant per
particle. From \eqref{eq:gas-final-rules},
\(\Im u_*=-(\pi/h_*)S_w\) and \(\Re u_*=O(N_0)\).
Combining the tilt with \eqref{eq:legacy-P-P3} and completing a square
in \({\cal D}+j\) therefore bounds the logarithm of the absolute
integrand by
\[
 -c\int({\cal D}+j)^2+C_{m,T',r}(N_0+1),
\]
locally uniformly in \(r\). Since \(|{\cal D}|\le2N_0\), the
\(j\)-sum of the first factor is at most \(C_{T'}(N_0+1)\).
The finite-volume factorial measures then sum the remaining exponential
in \(N_0\). Thus the final sum converges absolutely on a full imaginary
\(r\)-period, and the coefficient argument above identifies it with
\(P_{m,T'}(e^r)\).

\paragraph{Continuation in the complex period.}
For the size estimate, now set
\(a_0=\pi/b\), \(g=\sqrt{\log m}\), and
\(R=2a_0^{-1}\log m+g\). We describe continuation to
\(|T-R|<2\pi\), writing \(T'=\Re T\). Take all center contours in \eqref{eq:gas-grand-sum} (with offsets already absorbed into symbols) as the same period curve, a graph over the real period interval of length \(T'\) with Lipschitz constant of the imaginary displacement \(O(1/g)\), zero displacement except at \(|\Re s|>a_0^{-1}\log m+g/4\); impose this Lipschitz bound on the periodic translates taken together. This can be parametrized holomorphically in \(T\), for instance \(s=t+(T-R)\chi(t)\), \(\chi(t+R)=\chi(t)+1,\ \chi(\pm R/2)=\pm1/2,\ |\chi'|\le C/g\), with \(\chi=0\) throughout a suitable slightly wider central interval than specified. The curve is kept in a bounded imaginary strip on the central period. In \({\cal D},S_w,{\cal Q}\) now use the \textbf{real projections} of positions. Estimate \eqref{eq:legacy-P-P3} continues to hold for \(\Re{\cal E}_Z\) computed by analytic transport, with possibly smaller \(c_1\).

Here is the justification, which avoids using divergent oscillatory series on displaced curves. Decompose the raw symbol as \(\alpha WW^t/p^2+M_*e^{-p^2}/p+Y_1(p)\). The Laurent rule gives exactly \(T/12-\operatorname{sgn}(\Re d)d/2\) as the pair kernel for \(1/p^2\) (\(d=s-s'\) in the period, by the Fourier series for the quadratic Bernoulli polynomial, including the exponential's finite part). Thus its real part is unchanged from the comparison at \(T'\). The \(e^{-p^2}/p\) term is likewise the kernel \(-i\,\operatorname{sgn}(\Re d)/2\) plus periodized short smoothing correction: explicitly multiply \(-i/2\) by
\[
 \operatorname{sgn}(\Re d)+\sum_\nu
 [\operatorname{erf}((d+\nu T)/2)-\operatorname{sgn}(\Re(d+\nu T))].
\]
This follows by differentiating (the derivative uses the periodized Gaussian) and oddness. Here and below \(|\Im(d+\nu T)|\le Cg^{-1}|\Re(d+\nu T)|\), so displacement on this term costs \(O(g^{-1}{\cal Q})\).

Finally \(Y_1\) is regular, hence uses ordinary periodization of its inverse Fourier transform. Away from contact that transform has derivative \(O((1+1/|\Re s|)e^{-\gamma|\Re s|})\) continuing into sectors \(|\Im s|\le\theta|\Re s|\) for some \(\theta,\gamma>0\). Indeed for \(\Re s>0\) shift the inverse Fourier integral into the lower half-plane a small fixed distance near the origin, bending both tails further downward at a small fixed slope. The displayed rational functions in \(e^{Lp/12}\) have no intervening poles (take a sufficiently small shift near the axis), and \(Y_1=O((1+|p|)^{-1})\) on the tails. This gives the estimate, and similarly in the upper half-plane for \(\Re s<0\). Contour rotation agrees with inverse transformation on the real line by damping or truncation (for the undifferentiated integral the end segments vanish off the origin). The transform has just integrable logarithmic behavior there, given by the high-frequency terms. Consequently the displacement again costs \(O(g^{-1}{\cal Q})\) by the derivative bound on each translate. For self terms the image with no shift evaluates the regular diagonal independent of \(T\); the remaining translates obey the same estimate. These periodizations at real period follow equally by ordinary Fourier coefficients of the resulting locally integrable functions. Together they give the claimed analytic prescription and \eqref{eq:legacy-P-P3}.

For the linear term the final chemical rule gives per particle of type \(i\) exactly
\[
 i\,2\pi m\,\sigma_i\int_0^{s+i\delta_i L}
       \sum_{\nu\in\mathbb Z}\rho(d+\nu T)\,dd,\qquad
 \sigma=(0,1,1,-1,0).
\]
Indeed the primitive has this zero by oddness about 0, and uses all modes including the density's mean by the finite-part rule. Here
\[
 \rho(d)=\frac1{\sqrt3 L}\frac{\sinh(2a_0 d)}{\sinh(3a_0 d)}
\]
by summing the geometric progression of residues of \(\widehat\rho\) (poles at \((6k\pm1)i a_0\)). On \(|\Im d|\le 5L/12\) its real part is bounded below by \(c e^{-a_0|\Re d|}\): multiplying numerator by conjugate denominator shows strict positivity (both cosines positive and sines of the same sign), with the limits given by the formula. Each nonzero-\(\sigma\) linear real part per particle therefore satisfies
\[
 \le -c_2 m e^{-a_0|\Re s|}+O(1),\qquad c_2>0.
\]
In fact integrate along the real line then vertically, using \(\sigma_i\delta_i>0\) on those species. Contributions due to imaginary displacement of \(s\) itself, which is away from the core, and to images \(\nu\ne0\) throughout both integration legs, cost \(O(1)\) after multiplying by \(m\), by the exponential decay also along these deformations.

Write \(r/(2\pi)=u+i v\) with \(r\) in a compact set. From the final formula for \(u_*\), \(\Im u_*=-(\pi/h_*)S_w,\ \Re u_*=O(N_0)\). Combining \eqref{eq:legacy-P-P3}, the tilt and the \(T r_j^2\) term now bounds log absolute weights (excluding \({\cal P}\) and the factorial differentials) by
\[
 \frac{\pi^2}{h_*}
 \left[T'(u^2-v^2)-\int\big(\eta({\cal D}+j)^2+2v({\cal D}+j)\big)\right]
 + C(N_0+|j|+1) - c_2\sum_{\sigma\ne0}m e^{-a_0|\Re s|}.
\]
When \(2|v|<.54\), the integral dominates \(c\int({\cal D}+j)^2\) since \({\cal D}+j\) is integer valued. For general bounded \(r\) the same holds with loss \(O(T')\), sufficient for absolute convergence on a full period. The term \(O(|j|)\) is harmless since \(|j|\le 2N_0+(\int({\cal D}+j)^2/T')^{1/2}\).

\paragraph{Localization of the grand sum.}
\label{gas:grand-localization}
For localization put \(f_0(t)=(T'/2-g-|t|)_+\) periodically on the real projection. This is 1-Lipschitz, so \(|\sum W_d f_0(\Re s_d)|\le\int|{\cal D}+j|\le\int({\cal D}+j)^2\). The only negative \(W_d\) has \(\sigma\ne0\); on the support of \(f_0\) the function \(m e^{-a_0|t|}\) dominates \(f_0(t)\). Thus we retain suppression \(\exp[-c\sum f_0(\Re s_d)-c\int({\cal D}+j)^2+C N_0]\) for constants \(C<\infty,\ c>0\). The \(j\)-sum of the staircase suppression is bounded by completing a square, and the one-particle integral of \(e^{-c f_0}\) is \(O(g)\) including the bounded Jacobians. The factorial grand sum thus costs just \(\exp(O(g))\).

The localization estimates also justify analytic continuation: the indicated parametrizations, periodized pieces, and regular diagonal constants give holomorphic integrands in \(T\) off null coincidences, with uniformly summable majorants just established on compact subdisks. (For mere local uniform convergence at fixed \(m\), discard localization and use \(\int({\cal D}+j)^2\ge T'[(|j|-2N_0)_+]^2\) uniformly.) The constant phase conventions in the transform can all be kept fixed (branch factors from the logarithms/normal ordering, Weyl signs and contour/residue orientations, independent of \(T\)); equivalently continue each real-period factor on this parametrization. Thus \eqref{eq:gas-grand-sum} equals \(P_{m,T}(e^r)\) throughout the complex-period disk, by its identity on real \(T\).

\paragraph{The scalar factor and maximum modulus.}
It remains to track \({\cal P}\). The part multiplied by \(m^2\) in its log is the all-mode sample sum with \(1/T\) for the regular symbol
\begin{equation}\label{eq:gas-bulk-symbol}
 \widehat V(p)\widehat\rho(p),\qquad \widehat V(p)=\frac{2\pi}{p}\frac{\cosh(3Lp)-1}{\sinh(4Lp)} .
\end{equation}
Here \(\widehat V(p)\) denotes the transform of the function in the definition of \(I\); this follows also from the potential array, since its lifted log at offset 0 is precisely the periodization of \(V(s)\). Periodizing \(\widehat V(p)\widehat\rho(p)\) in dual space thus gives \(I\) for the unshifted term. For positive translates \(\nu T\) its inverse has expansion
\[
 A e^{-2\nu T}+O(e^{-a_0\nu T'})
\]
from the simple poles at \(-2i,-a_0 i\), and similarly at negative translates by evenness. This estimate and continuation follow by shifting slightly beyond these poles and bending tails as above. Therefore
\begin{equation}\label{eq:gas-bulk-period-correction}
 \frac1T\sum_p\widehat V(p)\widehat\rho(p)
 =I+2A e^{-2T}+O(e^{-a_0T'}),
\end{equation}
with \(m^2 e^{-a_0T'}=o(1)\) for the present growing buffer.
The log modulus from the determinant product is
\begin{equation}\label{eq:gas-determinant-rate}
 -\frac{T'}{\pi L}\int_0^\infty
 \log\big[(1-e^{-x}+e^{-2x})(1-e^{-6x})/(1-e^{-8x})\big]\,dx
 +o(T')=\frac{\pi T'}{16L}+o(T').
\end{equation}
Indeed these are converging Riemann sums (also for the periods with bounded imaginary parts, by analyticity near the ray and exponential decay). Expand the logs using \(1-t+t^2=(1+t^3)/(1+t)\) and \(\sum n^{-2}=\pi^2/6\).

Finally at fixed \(r\) multiply \(P_{m,T}(e^r)\) by
\(\exp(-2A m^2 e^{-2T})\) before using maximum modulus on
\(|e^{-2T}|\le e^{-2R}\). Analyticity near zero here has radius bounded below independently of large \(m\) by the original small residue cycles: take all centers close to the indicated pole with their internal residues already taken (no remaining pair poles there), so all product arguments stay bounded and translated image factors are analytic and nonzero. Its value at the center is exactly \(P_m(e^r)\). The estimates on the circle (using \(|\Im T|\le\pi/2\)) now give the upper bound. Indeed
\(a_0=8/3\), \(h_*=9\pi^2/32\), and
\(T'=R=(3/4)\log m+g\): the determinant contributes
\((3/8)\log m+o(\log m)\), the Gaussian tilt contributes
\(2\Re(r^2)\log m/(3\pi^2)+o(\log m)\), and the localized
grand sum costs only \(O(g)=o(\log m)\).
\end{proof}

The upper bound leaves the physical normalization undetermined. The next section supplies a saturation point; its nonvanishing will also be needed when Section~\ref{sec:uniform-nesting} sharpens the logarithmic error to a bounded one.
